\section*{Limitations}

\paragraph{One signer in the live test.} The live evaluation uses \V{lg.clips} clips recorded by a
single signer, the author (\authortodo{state your ASL experience}), in one room with one laptop
camera. If the signer's productions differ from the dataset's signers, some live errors are
production differences rather than pipeline or camera effects, and nothing in this design
separates them; a rating of each clip by a fluent, ideally Deaf, signer would. The results show
which pipeline discrepancies matter for this signer and camera; they do not estimate accuracy for
Deaf signers in general. With fifty clips most effects below the detection limit stated in
Section~\ref{sec:livegap} are inconclusive. The ten-clip pilot was recorded by the same person.

\paragraph{Controls we did not run.} Two controls would locate the remaining gap: running the live
pipeline on test-split landmark sequences, which measures what the live code costs in
distribution, and an oracle segmentation of the live clips, which separates segmentation from
recognition. The segmenter's thresholds rest on an assumed dataset frame rate. We also did not
disentangle architecture from model size or recipe in the v2 ablation.

\paragraph{Isolated signs from a closed vocabulary.} The model recognizes \V{data.signs} isolated
signs. It does not handle continuous signing, fingerspelling, or signs outside the vocabulary,
which it maps to the nearest known sign unless the confidence threshold rejects them. It is not a
translator: the gloss row and the generated sentence drop the grammar that ASL carries in space,
movement and facial expression.

\paragraph{Language and population.} The data is American Sign Language. The author is based in
India, where Indian Sign Language is the relevant language; nothing here transfers to it without
new data, such as INCLUDE \citep{sridhar2020include}. We know only the number of signers in the
dataset (\V{data.signers}), not their demographics, hearing status or fluency, so we cannot say
for whom the model works less well.

\paragraph{Test-set statistics and reuse.} With \V{data.testsigners} test signers, sequence-level
tests treat thousands of sequences from the same person as independent; the per-signer results
(Table~\ref{tab:persigner}) are the more conservative evidence, and a four-of-four split can never
be significant on its own. The test split was evaluated several times during development
(Table~\ref{tab:testaccess}), so its numbers are less clean than a single final evaluation.

\paragraph{Hardware and trackers.} Browser timings come from one laptop and one browser version;
the CPU fallback of MediaPipe reaches only \V{t.cpu.fps} frames per second there
(Section~\ref{sec:efficiency}).
